\section{Задание}
\subsection{Задание 1}
Создать базу знаний <<Предки>>, позволяющую наиболее эффективным способом 
(за меньшее количество шагов, что обеспечивается меньшим количеством предложений БЗ --- правил), 
и используя разные варианты (примеры) простого вопроса, определить:
\begin{itemize}
    \item по имени субъекта определить всех его бабушек;
    \item по имени субъекта определить всех его дедушек;
    \item по имени субъекта определить всех его бабушек и дедушек;
    \item по имени субъекта определить его бабушку по материнской линии;
    \item по имени субъекта определить его бабушку и дедушку по материнской линии.
\end{itemize}

Минимизировать количество правил и количество вариантов вопросов. 
Использовать конъюнктивные правила и простой вопрос. Для одного из 
вариантов вопроса задания 1 составить таблицу, отражающую конкретный 
порядок работы системы.

\subsection{Задание 2}
Дополнить базу знаний правилами, позволяющими найти:
\begin{itemize}
    \item максимум из двух чисел:
    \begin{itemize}
        \item без использования отсечения;
        \item с использованием отсечения;
    \end{itemize}
    \item максимум из трех чисел:
    \begin{itemize}
        \item без использования отсечения;
        \item с использованием отсечения;
    \end{itemize}
\end{itemize}

Убедиться в правильности результатов. Для каждого случая пункта 2 
обосновать необходимость всех условий тела. Для одного из вариантов 
вопроса и каждого варианта задания 2 составить таблицу, отражающую 
конкретный порядок работы системы.